% === Begin Label Data ===
% ID: PHY-JH-0004
% 学科: 物理
% 学段: 初中
% 年级: 九年级
% 册别: 通用
% 章节: 实验探究
% 知识点: 探究滑动摩擦力大小与接触面粗糙程度的关系
% 题型: 实验探究题
% 难度: 1.5
% 分值: 6
% 实验类型: 力学实验
% 图像类型: 无
% 单位要求: 力使用 $\mathrm{N}$
% 是否包含电路图: 否
% 是否包含受力图: 否
% 是否包含光路图: 否
% 是否包含实验表格: 否
% 来源: 本地原创测试题
% 是否启用: 是
% 备注: 初中实验探究基础测试题
% === End  Label Data ===

\begin{problem}
在“探究滑动摩擦力大小与哪些因素有关”的实验中，小组同学用弹簧测力计匀速拉动同一木块。

当木块放在木板表面时，弹簧测力计示数为 $1.2\,\mathrm{N}$；当木块放在毛巾表面时，弹簧测力计示数为 $2.0\,\mathrm{N}$。在其他条件相同的情况下，比较两次滑动摩擦力的大小，并写出实验结论。
\end{problem}

\begin{answer}
木块放在毛巾表面时受到的滑动摩擦力较大。

实验结论是在压力相同的情况下，接触面越粗糙，滑动摩擦力越大。
\end{answer}

\begin{solutions}
木块做匀速直线运动时，弹簧测力计的拉力大小等于滑动摩擦力大小。

因此在木板表面时，滑动摩擦力为 $1.2\,\mathrm{N}$；在毛巾表面时，滑动摩擦力为 $2.0\,\mathrm{N}$。

比较可知 $2.0\,\mathrm{N}>1.2\,\mathrm{N}$，所以毛巾表面的滑动摩擦力更大。

由此可得，在压力相同的情况下，接触面越粗糙，滑动摩擦力越大。
\end{solutions}
